\section{Caso de negocio}\label{b-0}

\subsection{La oportunidad}\label{b-2}

El gasto en la nube crece de forma sostenida, y la disciplina FinOps ---la práctica de dar visibilidad y control al coste cloud--- se ha convertido en una necesidad. Sin embargo, las herramientas FinOps del mercado están diseñadas para grandes organizaciones: asumen equipos dedicados, integraciones complejas y presupuestos que una startup o un profesional independiente no tienen. El resultado es una audiencia amplia y desatendida que paga de más porque no dispone de las herramientas ni del tiempo para entender su factura.

\subsection{Audiencia objetivo}\label{b-6}

\begin{itemize}
\item \textbf{Fundadores de startups} que gestionan su propio gasto cloud sin un equipo dedicado.
\item \textbf{Desarrolladores independientes} que pagan de su bolsillo la infraestructura de sus proyectos.
\item \textbf{Equipos pequeños} sin presupuesto para herramientas FinOps empresariales.
\end{itemize}

\subsection{Análisis de mercado}\label{b-23}

Gartner prevé que el gasto mundial de los usuarios finales en nube pública alcance los 723.400 millones de dólares en 2025, frente a los 595.700 millones de 2024 \cite{gartner2024}. Según Flexera, que encuestó en 2025 a 759 profesionales cloud, el 84\,\% sitúa la gestión del gasto entre sus principales retos y estima en un 27\,\% el gasto en infraestructura y plataformas en la nube (IaaS y PaaS) que se desperdicia \cite{flexera2025}.

Este problema ha impulsado el mercado de las herramientas FinOps. Sin embargo, el sector ha mantenido una orientación clara hacia las grandes empresas, pese a que el desperdicio es proporcionalmente mayor en las pequeñas, que son precisamente quienes menos medios tienen para corregirlo. Aquí radica la oportunidad de Economicon: startups, empresas pequeñas o profesionales independientes, donde el problema, en porcentaje, es más grave y las soluciones existentes no terminan de encajar.

Los datos según tamaño de organización confirman la paradoja: la tasa de desperdicio y el desperdicio en valor absoluto se mueven en sentidos opuestos a medida que la empresa crece. Ambos tipos de organizaciones tienen problemas distintos que requieren soluciones distintas. Economicon se centra en el primer grupo, el peor cubierto hoy.

El mercado, además, está saturado de herramientas de visualización dirigidas a quien ya sabe de FinOps: paneles que muestran el gasto y esperan que el usuario sepa interpretarlo. En el segmento pequeño no competimos con la solución empresarial, sino con la herramienta nativa y gratuita del propio proveedor, que casi nadie aprovecha porque no la entiende.

Ese hueco, el de la traducción y no el de los datos, es el que Economicon viene a ocupar.

\subsection{Justificación de la solución}\label{b-35}

\subsubsection{Diferenciador central}
Un agente conversacional que responde «¿en qué se me va el dinero?» con datos y fuentes, en lugar de exigir que el usuario construya la respuesta a partir de un panel.

La combinación de un motor analítico determinista, para los números (indicadores clave, anomalías y evaluación de presupuestos), con una capa conversacional basada en RAG, para la interpretación y la explicación, responde directamente al problema identificado. El motor analítico garantiza que las cifras son exactas y trazables; la capa conversacional elimina la barrera de conocimiento que impide al segmento objetivo aprovechar esas cifras. Se ha priorizado el vertical de Azure para el MVP por disponer de un conjunto de datos público representativo y de un modelo de datos de coste bien documentado, lo que permite validar el circuito completo sin comprometer datos reales de clientes.

\subsection{Viabilidad}\label{b-41}

\subsubsection{Técnica}
El prototipo implementa ingesta y normalización, indicadores, evaluación de presupuestos, monitorización y conversación con contexto documental. Hay pruebas automáticas y recorridos observados con datos reproducibles, sin una cuenta real de cliente. La aceptación integral del asistente sigue pendiente por los fallos y tiempos de respuesta detallados en el apartado~\ref{h-0}; implementar y probar el circuito no implica que todos sus resultados sean conformes. \cite{arquitectura-final,demo-final}

\subsubsection{Operativa}
En local, un único comando de Docker Compose, una herramienta que arranca varios servicios a la vez, levanta doce servicios con proveedores simulados, y catorce al activar la IA. Se alimenta de un conjunto de datos público de Microsoft (FinOps Toolkit v14, con licencia MIT) servido por una API de Azure simulada, de modo que no necesita conectarse a una cuenta real.

\subsubsection{Económica}
La tabla~\ref{tab-b1} compara dos escenarios económicos hipotéticos, sin una prueba de capacidad que acredite atender a ese número de clientes. Se suponen un pago único de despliegue de 200 € por cliente y una mensualidad de 50 €, con todos los clientes activos durante un año. Se presupuestan 1.200 € mensuales de infraestructura y 80 horas de mantenimiento y evolución a 25 € por hora: 3.200 € al mes. Mantener ese coste fijo en ambos escenarios es una hipótesis del cálculo; quedan por validar concurrencia, volumen de datos y coste por uso.

\begin{table}[htbp]
\centering\small
\caption{Escenarios económicos a un año: hipótesis internas, en euros}\label{tab-b1}
\begin{tabular}{@{}>{\raggedright\arraybackslash}p{0.46\linewidth}>{\raggedleft\arraybackslash}p{0.22\linewidth}>{\raggedleft\arraybackslash}p{0.22\linewidth}@{}}
\toprule
\textbf{Concepto} & \textbf{Conservador (20 clientes)} & \textbf{Optimista (100 clientes)} \\
\midrule
Ingresos por despliegue (200 € por cliente) & 4.000 & 20.000 \\
Ingresos por mensualidad (50 € durante 12 meses) & 12.000 & 60.000 \\
Costes (3.200 € al mes) & 38.400 & 38.400 \\
\textbf{Resultado} & \textbf{$-$22.400} & \textbf{41.600} \\
\bottomrule
\end{tabular}
\end{table}

Bajo esas hipótesis, el escenario conservador no cubre los costes y el optimista sí. El equilibrio recurrente requiere 64 clientes, o 48 en el primer año contando el pago de despliegue; estos umbrales no acreditan que la infraestructura soporte esa carga. Como hipótesis adicional de beneficio para el cliente, recuperar un tercio del desperdicio medio citado por Flexera supondría ahorrar el 9\,\% del gasto: una mensualidad de 50 € se compensaría a partir de unos 556 € al mes. Ese ahorro no se ha medido en usuarios de Economicon. \cite{flexera2025}

\subsection{Impacto esperado}

El impacto que persigue Economicon es que una persona sin formación en FinOps obtenga una respuesta trazable a la pregunta de en qué se le va el dinero, sin tener que construirla a partir de un panel. Para comprobarlo se midieron tres indicadores sobre una batería de 28 consultas de referencia: la proporción de respuestas correctas, con un umbral de 0,8; la latencia de respuesta en el percentil 95, es decir, el tiempo que no superan el 95\,\% de las respuestas, con un objetivo de 10 segundos; y la proporción de cifras de la respuesta que no se pueden rastrear hasta los datos.

La primera medición, hecha con una versión del chat que todavía no usaba un modelo y devolvía una plantilla con los fragmentos recuperados, dio 0 de 20 casos de respuesta aceptados \cite{jup070}.

\subsection{Costes y supuestos}

Los costes se presentan por separado: mediciones de IA con una población y un método definidos, y un presupuesto hipotético de infraestructura para el caso de negocio. Las ventanas de consumo no constituyen una factura acumulada del proyecto, y el presupuesto no mide capacidad ni fija un máximo de factura.

\subsubsection{Coste de la inteligencia artificial}
Una ventana histórica con DeepSeek V4.1 Flash registró unos 0,0006 USD por consulta terminada, unos 0,06 USD por cien, con contraste de facturación. No es el coste unitario de todas las campañas posteriores. \cite{jup133-coste} El ADR del gateway registró 0,0118 USD para tres tandas diagnósticas y una comprobación de embeddings de aquel benchmark; esa cifra no es el gasto acumulado del desarrollo. \cite{adr-registro} Las seis repeticiones recientes registraron 0,5291594448 USD por 171 operaciones previstas, incluidos reintentos y embeddings, según el gateway y sin contraste de factura. Se presentan con su denominador y límites en el apartado~\ref{h-0}; no se suman a otras ventanas que puedan solaparse. El coste del juez IA usado para revisarlas no está medido en ese registro. \cite{campanas-infra,juez-infra}

\subsubsection{Coste de infraestructura}
Compose define diez servicios permanentes y dos de inicialización; el perfil de IA añade dos servicios. La referencia AWS usa servicios gestionados y no acredita un despliegue de la aplicación. \cite{arquitectura-final,aws-ensayo} Para el caso de negocio se usa un presupuesto supuesto de 1.200 € mensuales, sin IVA ni capa gratuita y con operación continua. La capacidad para cada escenario de clientes y la factura efectiva quedan por medir. La IA es un coste variable separado: los modelos nativos de Amazon y la ruta de Flash tienen precios y consumos distintos, de modo que una estimación de una ruta no representa la otra.

\subsection{Riesgos}

Los riesgos del proyecto se agrupan en tres bloques: los técnicos, los económicos y los de contexto.

\subsubsection{Riesgos técnicos}
El principal fue que el asistente no alcanzara la calidad necesaria. La línea base sin modelo, descrita en el impacto esperado, dio el punto de partida: ninguno de los 20 casos evaluados fue aceptado \cite{jup070}. La probabilidad y el impacto eran altos, y se controló midiendo siempre sobre la misma batería de consultas de referencia y con los mismos indicadores, de modo que la mejora pudiera demostrarse y no solo afirmarse. Ligado a este estaba el riesgo de latencia: las mediciones de agosto, de 11,73 s y 10,88 s, quedaban por encima del objetivo de 10 s. Su probabilidad e impacto eran medios, y se acotó con límites de unidades de texto y un tiempo de espera de 30 s.

\subsubsection{Riesgos económicos}
El uso de modelos de pago puede generar gasto no previsto por cambios de precio, reintentos o mayor uso. Se controla con límites por clave y registro de consumo; las mediciones históricas pequeñas no garantizan el coste futuro. \cite{adr-registro,campanas-infra} La infraestructura añade un coste dependiente de servicios, volumen y duración: el escenario económico se deberá revisar con carga y facturación observadas antes de ofrecer una capacidad o un precio garantizados.

\subsubsection{Riesgos de contexto}
El segmento objetivo, las organizaciones pequeñas, es una hipótesis de partida que no se contrastó con usuarios. Su probabilidad era media y su impacto alto, y se controló presentándolo siempre como hipótesis y dejando el contraste con usuarios como trabajo futuro. Por último, el plazo: la congelación de la memoria estaba prevista para el 16 de octubre y la entrega para el 23, sujeta a confirmación institucional. La probabilidad era media y el impacto alto, y se controló priorizando el contenido obligatorio, trabajando en pequeñas entregas revisables por el equipo y señalando de forma explícita lo que quedó pendiente.
